\documentclass[11pt]{article}
\usepackage[a4paper,margin=25mm]{geometry}
\usepackage{amsmath,amssymb,amsthm,mathtools}
\usepackage[T1]{fontenc}
\usepackage[hidelinks]{hyperref}
\newtheorem{theorem}{Theorem}
\newtheorem{lemma}[theorem]{Lemma}
\newtheorem{proposition}[theorem]{Proposition}
\newtheorem{corollary}[theorem]{Corollary}
\newcommand{\C}{\mathbb C}
\newcommand{\R}{\mathbb R}
\newcommand{\diag}{\operatorname{diag}}
\newcommand{\adj}{\operatorname{adj}}
\newcommand{\Res}{\operatorname{Res}}
\newcommand{\Disc}{\operatorname{Disc}}
\newcommand{\ord}{\operatorname{ord}}
\title{Complete lifted actions and a real-time mixed boundary\\in the original ES--Fable conductor correspondence}
\author{The Clankers: additive mathematical continuation}
\date{20 September 2026}
\begin{document}
\maketitle
\begin{abstract}
We classify every constant complex four-by-four target matrix whose unique
polynomial lift through the original ES--Fable map is complete. Real-time
completeness for all initial points, complex-time completeness, and membership
in one specified weight line are equivalent. This strengthens the preceding
infinitesimal-equivariance test: a nonlinear lift cannot evade the classification.
For the literal marked arithmetic quartet we then use its full conjugated
matrix, rather than assigning its eigenvalues to the four polynomial coefficient
slots. An explicit real-time target orbit reaches a mixed leading-coefficient
and discriminant boundary. Its eight inverse states consist of three finite
limits, four quarter-power escapes and one cubic-pole escape. Their original
Gamma-norm constants and both directional singular spectra are evaluated.
The four root labels and both factor signs remain. The completed factor space
repairs the one chart escape, while the four double-root escapes persist.
The unchanged original weighted conductor transports every branch and remains
invertible. No asymptotic sign of the programme's independent projected
arithmetic cross product is inferred.
\end{abstract}

\section{Exact source, labels and scope}
The frozen provider is the complete signed-state edition \cite{Signed}, at
repository commit \texttt{d73b545c7e70dd55fb8b58557f704666ba310626}, blob
\texttt{50b39a82568bbe5104aeed3de5ad9441cb6a5602}. Its FC31 and FC44 supply
the original polynomial and factor chart. This resolves the coordinate-identity
qualification in the earlier reconstructed-map note: the following factor
formulas expand coefficient for coefficient to FC31. The original four-dimensional
polynomial and its four distinguished representatives are attributed by the
provider to the canonical version-69 Erd\H{o}s--Straus reader, theorem
\texttt{explicit-four-time-Jacobian-collision}, source SHA-256
\texttt{678b3f70ee7d2a6d00f6c588d7778e390118ab4743a5759612dbba6430984c78}.
We read the published provider, not a newly extracted copy of that archive.
Tao's exposition \cite{Tao} supplies the original linear--quadratic/resultant
background and attribution to Alp\"oge and Fable; the four-dimensional ES
extension and the calculations below have their separate provenance.

Write the source as $q=(a,y,z,w)\in\C^4$. The coordinate $a$ is distinct
from the programme's fixed amplification integer $a_{\rm amp}$. Put
\begin{align}
 b&=i+ay,&c&=-i+2ay+a^2z,\nonumber\\
 d&=-iy-a(iz+2y^2)-a^2yz,&e&=2z-7iy^2+aw,\nonumber\\
 f&=iw+3iy^3-4yz+a(6iy^2z+wy+4y^4)+2a^2y^3z.\tag{FRF1}
\end{align}
For $L=aU+bV$ and $Q=cU^3+dU^2V+eUV^2+fV^3$, direct substitution gives
\[
 ad+bc=1,\qquad \Res(L,Q)=-cb^3+dab^2-ea^2b+fa^3=1.
 \tag{FRF2}
\]
The original map is
\[
 P(q)=(ac,ae+bd,af+be,bf)=(A,B,C,D),\qquad
 F_Y(u)=Au^4+u^3+Bu^2+Cu+D.\tag{FRF3}
\]
Thus $LQ$ is the homogeneous quartic associated with $F_Y$. Its $U^3V$
coefficient is exactly one. On $a\ne0$ its inverse factor coordinates are
\[
 y=(b-i)/a,\quad z=(c+i-2ay)/a^2,\quad
 w=(e-2z+7iy^2)/a.\tag{FRF4}
\]
The factor chart at $a=0$ contains $b=i,c=-i$ and excludes $b=-i,c=i$.
The second chart is obtained by replacing $i$ with $-i$ in FRF1.

For reference, the four original marked source columns are
\[
 K=\begin{pmatrix}
0&i&1/\sqrt2&1/\sqrt2\\
0&-1&-1-\sqrt2i&1-\sqrt2i\\
i/2&-3i&2\sqrt2+6i&-2\sqrt2+6i\\
0&13&-34-19\sqrt2i&34-19\sqrt2i
\end{pmatrix},\qquad \det K=77\sqrt2i/2.\tag{FRF5}
\]
Its columns have labels $(\mathrm M,\mathrm N,\mathrm T,\mathrm E)$.
They map to $(0,0,-1,0)$; their selected projective roots are, respectively,
$\infty,0,1,-1$. The sign state represented by each column is denoted $+$;
its opposite factorization has sign $-$. These are eight factor states,
with the $\mathrm M,-$ state outside the positive chart at that target.
The informal earlier label $\mathrm A$ is not identified with $\mathrm N$.

\begin{lemma}[Original Jacobian and all finite fibres]\label{lem:fibres}
The map in FRF3 satisfies $\det DP=-2$ everywhere. For every target $Y$,
\[
 \#P^{-1}(Y)=2n_s(Y)+\mathbf1_{\{A=0\}},\tag{FRF6}
\]
where $n_s(Y)$ is the number of finite simple roots of $F_Y$.
For a finite simple selected root $r$, the two factor signs satisfy
\[
 a^2F_Y'(r)=1,\quad b=-ar,\quad
 (c,d,e,f)=a^{-1}(A,1+Ar,B+r+Ar^2,C+Br+r^2+Ar^3).\tag{FRF7}
\]
\end{lemma}
\begin{proof}
For $a\ne0$, the intermediate coordinates $(a,b,c,e)$ have determinant
$a^4$ relative to $(a,y,z,w)$. Eliminating $d,f$ in FRF2 gives
\[
 (A,B,C,D)=\left(ac,\ ae+b/a-b^2c/a,\
 (1-b^2+2cb^3)/a^2+2be,\
 b(1-b^2+2cb^3)/a^3+b^2e/a\right).
\]
Their determinant relative to $(a,b,c,e)$ is $-2/a^4$ by differentiation.
The chain rule and polynomial continuation prove the Jacobian assertion.
Polynomial division by $u-r$ gives FRF7, and evaluation of the resultant
gives $a^2F_Y'(r)=1$. A repeated selected root has resultant zero and is
excluded by FRF2. For $A=0$, infinity is a simple projective root because
the next coefficient is one. Its normalized factors have $b=\pm i$;
exactly one lies in the positive chart. FRF4 and its regular inverse at
$a=0$ prove FRF6. These are the original FC44/GF inverse formulas,
rederived here to specify the completeness argument's exact input.
\end{proof}

\section{Every complete polynomial lift of a linear target action}
For $M\in\operatorname{Mat}_4(\C)$ define
\[
 \mathcal X_M(q)=(DP(q))^{-1}MP(q)
       =-\tfrac12\adj(DP(q))MP(q).\tag{FRF8}
\]
This is a polynomial vector field. A lift of the target infinitesimal
linear action is unique because $DP(q)$ is invertible at every finite
point. The time variable below is independent of the original period,
tensor degree and real source tilt. Real-time completeness means existence
for every real time and every initial point in the real eight-dimensional
manifold underlying $\C^4$.

\begin{theorem}[Complete classification, FRF-C]\label{thm:complete}
For the fixed original $P$, the following three conditions on $M$ are equivalent:
its polynomial lift $\mathcal X_M$ is real-time complete; it is complex-time
complete; and
\[
 M=\tau\diag(1,-1,-2,-3),\qquad \tau\in\C.\tag{FRF9}
\]
On this line $\mathcal X_M(q)=Mq$, with its ordinary exponential flow.
This classifies nonlinear lifts, not only linear source actions satisfying
$DP(q)Mq=MP(q)$.
\end{theorem}
\begin{proof}
Let $\Delta(A,B,C,D)=\Disc_uF_Y$. If the source flow $\varphi_t$ is complete
for all real $t$, differentiation gives
$P\varphi_t=e^{tM}P$, and $\varphi_{-t}$ supplies its inverse. It follows
that the cardinality of every fibre is preserved by $e^{tM}$, including
zero-cardinality fibres. By FRF6 the odd-cardinality locus is exactly
$\{A=0\}$, so this hyperplane is preserved. Outside it, a fibre has eight
points exactly when $\Delta\ne0$; on it, a fibre has seven points exactly
when $\Delta\ne0$. Therefore the discriminant hypersurface is preserved
as well. In particular the first row of $M$ has the form $(\alpha,0,0,0)$.

We record why the reduced polynomial $\Delta$ is irreducible on this
literal coefficient slice. Every multiple projective root is finite,
because infinity is never multiple when the $U^3V$ coefficient is one.
The repeated-root incidence is parametrized surjectively by
\[
 (A,B,r)\longmapsto
 (A,B,-4Ar^3-3r^2-2Br,\ 3Ar^4+2r^3+Br^2).\tag{FRF10}
\]
Its image is precisely $\{\Delta=0\}$, by $F(r)=F'(r)=0$, so that
hypersurface is irreducible. At $(0,0,-3,2)$ the gradient of $\Delta$ is
$-108(1,1,1,1)$, as computed below; hence the displayed discriminant is
reduced. The principal ideal of the hypersurface is consequently $(\Delta)$.
A linear vector field preserving it must obey
\[
 (MY)\cdot\nabla\Delta=\lambda\Delta\tag{FRF11}
\]
for a constant $\lambda$: the left side has degree at most $\deg\Delta=6$.

At $A=0$ the discriminant is
\[
 \Delta_0=B^2C^2-4C^3-4B^3D-27D^2+18BCD.\tag{FRF12}
\]
The degree-two term $-27D^2$ in FRF11 first forces the $B,C$ coefficients
in the $D$ component of the vector field to be zero, and gives
$\lambda=2m_{DD}$. The degree-three terms then force the off-diagonal
coefficients in the $B,C$ components to vanish. Comparing $C^3$ and $BCD$
gives $m_{CC}=2m_{DD}/3$ and $m_{BB}=m_{DD}/3$.
Thus, writing $\beta=m_{BB}$, the full field has the form
\[
 \alpha A\partial_A+(\beta B+pA)\partial_B
 +(2\beta C+qA)\partial_C+(3\beta D+rA)\partial_D,
 \qquad \lambda=6\beta.\tag{FRF13}
\]
Here the coefficient $r$ in the vector field is separate from the selected
root in FRF10.

The coefficient of $A$ in the full discriminant is
\[
 \Delta_1=144BD^2-6C^2D-80B^2CD+18BC^3+16B^4D-4B^3C^2.
 \tag{FRF14}
\]
It has weight seven when $B,C,D$ have weights $1,2,3$. The $A$ coefficient
in FRF11 therefore gives
\[
 (\alpha+\beta)\Delta_1+
 p\partial_B\Delta_0+q\partial_C\Delta_0+r\partial_D\Delta_0=0.
\]
The degree-five terms force $\alpha=-\beta$, since the three derivatives
have degree at most three. The remaining degree-one term forces $r=0$;
the degree-two $C^2$ and $CD$ terms then force $q=p=0$. This proves the
necessary weight line. Conversely every monomial in FRF1--3 has the
claimed weights, so $DP(q)Mq=MP(q)$ on FRF9. Uniqueness in FRF8 gives
$\mathcal X_M=Mq$ and the entire exponential flow. Complex-time
completeness implies real-time completeness, completing the equivalence.
\end{proof}

\begin{corollary}[Original translation and quartet actions]
Every nonzero nilpotent target matrix has an incomplete lifted vector field.
No target matrix whose spectrum is the original offcritical quartet
$c\pm d\pm ig$, with $d,g>0$, has a complete lift, in any constant
coordinate conjugation.
\end{corollary}
\begin{proof}
A nilpotent matrix on the line FRF9 is zero. For the quartet, its trace
is $4c$ while the trace of FRF9 is $-5\tau$. If $c$ is real, this makes
$\tau$ real, whereas the quartet has nonreal eigenvalues. At $c=0$ it
instead forces $\tau=0$. Similarity preserves these conclusions.
\end{proof}

\section{The exact return to the original conductor and its actions}
Keep the original period and branch, full conductor coefficients $a_{uw}$,
actual vanishing order $v$, and first nonzero moment $\mu_v$. The quartet
in the fixed FC5 order is
\[
 R=\diag(\rho_{\mathrm M},\rho_{\mathrm N},\rho_{\mathrm T},\rho_{\mathrm E}),
 \quad \rho=(c+d+ig,c+d-ig,c-d+ig,c-d-ig).\tag{FRF15}
\]
For the literal one-packet application $c=1/2,d=\delta,g=\gamma$ with
$0<\delta<1/2,\gamma>2$. Write
$h(S)=\prod_\alpha(S-\rho_\alpha)$ and
$\ell_\alpha(S)=\prod_{\beta\ne\alpha}(S-\rho_\beta)/(\rho_\alpha-\rho_\beta)$.
The original maps FC26--35 are
\[
 Xe_\alpha=[S^v\ell_\alpha(S)]\in
 \mathcal U=\mathcal P_{v+3}/\mathcal P_{v-1},\qquad
 Ye_\alpha=T_A Xe_\alpha\in\mathcal W=\mathcal P_3,\tag{FRF16}
\]
\[
 T_AP(S')=\sum_{u,w=0}^8a_{uw}P(S'+\beta_{8,u,w}),\quad
 \Phi=XK^{-1},\quad\Psi=YK^{-1},\quad T_A\Phi=\Psi.\tag{FRF17}
\]
The letter $Y$ in FRF16 denotes a fixed matrix, distinguished from a
variable target coefficient vector by its occurrence as a linear map.
Every coefficient of the fixed matrix is explicit:
\[
 Y_{j\alpha}=\sum_{n=j}^3\lambda_{n\alpha}
       \binom{v+n}{j}\mu_{v+n-j},\qquad
 \ell_\alpha(S)=\sum_{n=0}^3\lambda_{n\alpha}S^n.\tag{FRF18}
\]
The determinant is
\[
 \det Y=\frac{\mu_v^4\prod_{j=0}^3\binom{v+j}{v}}
 {-64d^2g^2(d^2+g^2)}\ne0.\tag{FRF19}
\]
This follows from the triangular Taylor matrix in the shifted monomials
and the literal Vandermonde evaluation matrix. In particular $\Phi,\Psi$
are fixed isomorphisms, and
\[
 P_{\mathcal U}=\Phi P\Phi^{-1},\quad
 P_{\mathcal W}=\Psi P\Psi^{-1},\quad
 T_AP_{\mathcal U}=P_{\mathcal W}T_A.\tag{FRF20}
\]
We use these published receiving maps, rather than the arbitrary
four-corner injection in the earlier reconstruction note.

Multiplication by $S$ on $\C[S]/(h)$ is $R$ in the labelled cardinal
basis. Transporting that specific action to $\mathcal W$ gives
$\mathcal B_{\mathcal W}=YRY^{-1}$. Its matrix in the polynomial
coefficient coordinates of FRF3 is therefore
\[
 \boxed{M=\Psi^{-1}\mathcal B_{\mathcal W}\Psi=KRK^{-1}.}\tag{FRF21}
\]
It is not $R$ acting directly on $(A,B,C,D)$. This is an explicit transport
of the arithmetic quotient action; multiplication by $S'$ on untruncated
polynomials is a different map. The original full Taylor unit is retained
as its separate diagonal multiplication by $u_\alpha=(2\xi/h)(\rho_\alpha)$.
It commutes with $R$. When appended to both marked receiving maps, its
inverse and its factors in the pulled-back Gram remain; the computation
FRF21 is unchanged. No unit is assigned the value one.

The other original action, differentiation, is treated without such an
identification. If $V_\rho$ is the four-by-four Vandermonde and
\[
 N_v=\begin{pmatrix}0&v+1&0&0\\0&0&v+2&0\\0&0&0&v+3\\0&0&0&0\end{pmatrix},
\]
FC50--53 gives its Fable-coordinate matrix
\[
 M_{\rm tr}=K V_\rho^T N_v(V_\rho^T)^{-1}K^{-1},\qquad
 M_{\rm tr}^4=0,\quad M_{\rm tr}^3\ne0.\tag{FRF22}
\]
It retains the original quotient correction
$f' +v(f-f(0))/S$. Theorem~\ref{thm:complete} proves that the unique
nonlinear lift of these original translations is incomplete. This is
stronger than the earlier calculation that $P$ does not commute with
their uncorrected linear source flow.

For any $M$, the transported lifted fields satisfy the exact derivative
square
\[
 D P_{\mathcal W}\,(\Psi\mathcal X_M\Psi^{-1})
       =(\Psi M\Psi^{-1})P_{\mathcal W},\qquad
 T_A(\Phi\mathcal X_M\Phi^{-1})
       =(\Psi\mathcal X_M\Psi^{-1})T_A.\tag{FRF23}
\]
The latter equation concerns evaluation of vector fields under a fixed
linear isomorphism. It changes none of the coefficients of $T_A$.

\section{An explicit real-time orbit of the literal marked quartet}
Put $Y_0=(0,0,-3,2)^T$, so $F_{Y_0}(u)=(u-1)^2(u+2)$. Define
\[
 k_w=K^{-1}e_4=
 \left(\frac6{77},\frac1{77},\frac{-3-\sqrt2i}{154},
                         \frac{3-\sqrt2i}{154}\right)^T,\quad
 K^{-1}e_3=-2ie_{\mathrm M},\tag{FRF24}
\]
\[
 \eta=K^{-1}Y_0=6ie_{\mathrm M}+2k_w,\qquad
 Z(\theta)=K\diag(e^{\rho_\alpha\theta})\eta=e^{\theta M}Y_0.
 \tag{FRF25}
\]
Both inverse-column identities follow by multiplication by FRF5. Thus
FRF25 specifies every target coordinate by four original exponentials.
Here $\theta$ is a flow displacement only, not a source tilt.

\begin{proposition}[Two evaluated transverse derivatives]\label{prop:transverse}
The derivatives of the leading coefficient and of the polynomial at its
double root are, respectively,
\[
 \kappa=Z_A'(0)=\frac2{77}
       \left(g+2id-\frac{3ig}{\sqrt2}\right),\tag{FRF26}
\]
\[
 \Omega=\left.\frac{d}{d\theta}F_{Z(\theta)}(1)\right|_0
 =-c-\frac{337}{77}d+\frac{34-190\sqrt2}{77}g
       +\frac{i}{77}\left(4d-(255+3\sqrt2)g\right).\tag{FRF27}
\]
For the original quartet both are nonzero. Moreover
\[
 \Delta(Y_0)=0,\quad \nabla\Delta(Y_0)=-108(1,1,1,1),\quad
 \left.\frac{d}{d\theta}\Delta(Z(\theta))\right|_0=-108\Omega.
 \tag{FRF28}
\]
\end{proposition}
\begin{proof}
Differentiate FRF25 and multiply $KR\eta$. For a shorter exact expansion,
$M e_3=\rho_{\mathrm M}e_3$ and
\[
 (M e_4)_1=\frac{g+2id-3ig/\sqrt2}{77},\qquad
 \mathbf1^TM e_4=c+\frac{-53d+(17-95\sqrt2)g}{77}
                 +\frac{i[2d-(12+3/\sqrt2)g]}{77}.
\]
Since $Y_0=-3e_3+2e_4$, these give FRF26--27. The real part of $\kappa$
is $2g/77>0$. The real part of $\Omega$ is strictly negative for $c,d,g>0$;
its displayed imaginary part is also negative on the original
$0<d<1/2,g>2$ domain. FRF12, FRF14 and the remaining discriminant terms
give the gradient in FRF28. The chain rule completes the proof.
\end{proof}

For full finite control, the entire discriminant in these coordinates is
\begin{align*}
\Delta={}&256A^3D^3-192A^2CD^2-128A^2B^2D^2+144A^2BC^2D-27A^2C^4\\
 &+144ABD^2-6AC^2D-80AB^2CD+18ABC^3+16AB^4D-4AB^3C^2\\
 &-27D^2+18BCD-4C^3-4B^3D+B^2C^2.
\end{align*}
The sum of absolute coefficients is $1069$ and its total degree is six.
Let $m=\|M\|_\infty$ be the maximum absolute row sum and set
\[
 \boxed{T=\min\left\{
 \frac1{300(1+m)},\ \frac{|\kappa|}{9(1+m)^2},\
 \frac{54|\Omega|}{36\cdot1069\cdot9^6(1+m)^2}
 \right\}>0.}\tag{FRF29}
\]
This is a specified positive real number in the original quartet parameters.
Our real-time orbit, for $0\le\tau\le T$, is
\[
 Y(\tau)=Z(\tau-T),\qquad Y(T)=Y_0.\tag{FRF30}
\]
No change of the original period is involved.

\begin{lemma}[No earlier boundary is discarded]\label{lem:guard}
For $0\le\tau<T$, $A(\tau)\ne0$ and $\Delta(Y(\tau))\ne0$.
Two roots remain in $|u-1|<1/2$, one remains in $|u+2|<1/2$, and the
fourth tends to infinity as $\tau\uparrow T$.
\end{lemma}
\begin{proof}
Write $\theta=\tau-T$. Since $\|Y_0\|_\infty=3$ and
$m|\theta|\le1/300$, the orbit and its first two derivatives are bounded
coordinatewise by $9,9m,9m^2$. Thus
\[
 |A(\theta)-\kappa\theta|\le\tfrac92m^2|\theta|^2.
\]
For any monomial of degree at most six, twice differentiating its product
along the orbit gives at most $36m^2 9^6$ in absolute value. Therefore
\[
 \left|\Delta(Z(\theta))+108\Omega\theta\right|
 \le\tfrac12(36\cdot1069\cdot9^6)m^2|\theta|^2.
\]
FRF29 makes both linear terms dominate on $-T\le\theta<0$.

Also $\|Z(\theta)-Y_0\|_\infty\le9m|\theta|$. On the circle
$|u-1|=1/2$, the unperturbed polynomial has modulus at least $5/8$;
the perturbation is at most $(157/16)9mT<5/16$.
On $|u+2|=1/2$ the corresponding bounds are $25/8$ and
$(781/16)9mT<25/16$. The strict inequalities follow from
$mT\le1/300$. Rouch\'e's theorem, equivalently the argument principle
under this nonvanishing boundary homotopy, gives the stated two and one
roots. The remaining projective root tends to the simple infinity root
of the limiting binary quartic. There are no repeated roots before $T$
by the discriminant bound.
\end{proof}

\section{All eight original-coordinate inverse trajectories}
Every root of $F_{Y(\tau)}$ is simple for $\tau<T$. Choose its continued
label and each of the two continuous square roots in FRF7. The full source
coordinates, with no discarded coefficient, are
\[
 \boxed{q(r,a;A,B)=\left(
 a,\ -r-i/a,\ A/a^3+2r/a+3i/a^2,\
 7ir^2/a+(B-17r+Ar^2)/a^2-13i/a^3-2A/a^4
 \right),}\tag{FRF31}
\]
where $a^2F_Y'(r)=1$. Substitution into $P$ gives $Y$.
Equivalently the two remaining target coordinates are
\[
 C=a^{-2}-4Ar^3-3r^2-2Br,\quad
 D=3Ar^4+2r^3+Br^2-r/a^2.\tag{FRF32}
\]
These formulas specify eight initial points at $\tau=0$ and their full
continuation. Differentiating $P(q(\tau))=Y(\tau)$ proves
$q'=\mathcal X_M(q)$ exactly. There is no independent trajectory chosen
in place of the given action.

The four labels can be retained without guessing which cluster they
represent. Begin at the original product $u^3-u$ with the FRF5 marking,
and follow the cubic path
\[
 u^3-(1+2s)u+2s=(u-1)(u^2+u-2s),\quad 0\le s<1.\tag{FRF33}
\]
The $\mathrm T$ root stays at one, the $\mathrm N$ root
$(-1+\sqrt{1+8s})/2$ approaches one, the $\mathrm E$ root approaches
$-2$, and the $\mathrm M$ root stays at infinity. Close to $s=1$,
connect to $Y(0)$ inside the three disjoint projective root neighborhoods
of Lemma~\ref{lem:guard}, avoiding the discriminant. Such a connection
exists by moving the center and nonzero difference of the two roots in
the first disk, and the separate simple roots in the other disks. The
$U^3V$ coefficient is nonzero in this neighborhood and division by that
coefficient returns exactly to the original coefficient slice. The path
and both signs are retained. Different choices may exchange the two
labels inside the first disk or their signs, but do not change its set
$\{\mathrm N,\mathrm T\}$ or any count or rate below. Working in the
completed factor space retains the initially excluded $\mathrm M,-$
state during this labelling path.

\begin{theorem}[Real-time mixed-boundary census, FRF-B]\label{thm:branches}
Along FRF30 there are precisely three finite source limits,
\[
 (0,0,3i/2,-2),\qquad
 (1/3,2-3i,-12+27i,306-267i),\qquad
 (-1/3,2+3i,12+27i,306+267i).\tag{FRF34}
\]
They are the complete fibre $P^{-1}(Y_0)$. The other five trajectories
escape: the four states over the colliding $\mathrm N,\mathrm T$ roots
have norm growth $(T-\tau)^{-1/4}$, and the excluded infinity sign has
norm growth $(T-\tau)^{-3}$.

For $G=\Psi^*\mathcal H\Psi$ equal to the original target Gamma Gram
in the unchanged Fable coordinates, put $\varepsilon=T-\tau$. On each
of the four colliding states,
\[
 \boxed{\lim_{\tau\uparrow T}\varepsilon^{1/4}
       \|\Psi q(\tau)\|_{\rm Gamma}
       =\frac{\sqrt{G_{11}}}{\sqrt2(3|\Omega|)^{1/4}}>0.}\tag{FRF35}
\]
On the excluded infinity sign,
\[
 \boxed{\lim_{\tau\uparrow T}\varepsilon^3
       \|\Psi q(\tau)\|_{\rm Gamma}
       =\frac{40\sqrt{G_{44}}}{|\kappa|^3}>0.}\tag{FRF36}
\]
The original source-quotient versions hold with
$G=(\Phi)^*G_{\mathcal U}\Phi$ instead.
\end{theorem}
\begin{proof}
Near $u=1$ put $u=1+\xi$. The exact Taylor expansion at $(Y_0,1)$ has
\[
 F_{Z(\theta)}(1+\xi)=3\xi^2+\xi^3+\Omega\theta
                         +O(\theta\xi+\theta^2).
\]
There are two analytic Puiseux branches
$\xi=\pm\sqrt{-\Omega/3}\,\theta^{1/2}+O(\theta)$.
For example, substitute $\theta=s^2$, $\xi=s v$; after dividing by
$s^2$, the two simple roots of $3v^2+\Omega=0$ continue by the ordinary
implicit function theorem. Thus
\[
 F_{Z(\theta)}'(r)=\pm2\sqrt{-3\Omega}\,\theta^{1/2}+O(\theta),\quad
 |a|\sim\frac{|\theta|^{-1/4}}{\sqrt2(3|\Omega|)^{1/4}}.\tag{FRF37}
\]
Each root has two factor signs. In FRF31 their other coordinates obey
$y\to-1,z\to0,w\to0$, so $q/a\to e_1$. This proves four genuine
factor escapes and FRF35. It also proves the exponent without assuming
an asymptotic for the entire source Gram.

For the infinity branch, use the negative factor chart, which has a
finite limiting point $(a,y_-,z_-,w_-)=(0,0,-3i/2,-2)$ over $Y_0$.
The exact transition back to the positive chart is
\[
 (a,y_-,z_-,w_-)\mapsto
 (a,y_--2i/a,z_-+6i/a^2,
 w_-+14iy_-^2/a+28y_-/a^2-40i/a^3).\tag{FRF38}
\]
In its factor coordinates $c\to i$, so $A=ac$ gives
$a/A\to-i$. Since $A(\theta)/\theta\to\kappa\ne0$ and the negative
chart coordinates are bounded, the leading positive-chart term is
\[
 q(\theta)\sim -40 A(\theta)^{-3}e_4
\]
in the sense that $A^3q\to-40e_4$. This proves FRF36. The other infinity
sign remains finite. The two signs at the simple root $-2$ have
$a=\pm1/3$ because $F_{Y_0}'(-2)=9$. FRF31 evaluates their limits
as FRF34; the positive infinity chart gives the first point there.
Lemma~\ref{lem:fibres} gives the exhaustive count three at $Y_0$ and eight
before it. The fixed invertible maps $\Phi,\Psi$ preserve the asserted
convergence and escape, and their actual Gram forms give the constants.
\end{proof}

The completed factor variety
\[
 \mathfrak X=\{ad+bc=1,\ -cb^3+dab^2-ea^2b+fa^3=1\}\subset\C^6
\]
has the additional finite state $(0,-i,i,0,-3i,2i)$ at this target.
Thus its limiting fibre has four points, with the other four trajectories
still escaping. Adding the missing chart repairs exactly the cubic-pole
trajectory. It does not repair the real-time incompleteness of the four
colliding-root lifts: their factor coefficient $a$ itself diverges.
This separates chart completion from continuation across the resultant
boundary while retaining every state.

\section{Both directional singular spectra in the original metric}
The following local calculation gives a general finite-root boundary
result, beyond one specially selected escaping curve.

\begin{lemma}[Finite-root differential spectrum]\label{lem:spectrum}
In FRF31 let $|a|\to\infty$ while $A,B,r$ remain bounded and
$12Ar^2+6r+2B$ has a nonzero limit. In any fixed positive source and
target Hermitian metrics, the ordered singular values of $DP(q)$ have
scales
\[
 |a|^3,\quad |a|^2,\quad 1,\quad |a|^{-5}.\tag{FRF39}
\]
Here a scale means two-sided comparison by fixed strictly positive constants.
\end{lemma}
\begin{proof}
Write $J=DP(q)$ and retain its original column order $(a,y,z,w)$. Put
\[
 v_r=(1,-r^2,-2r^3,2r^4)^T,\quad
 v_{1,r}=(0,2,-4r,2r^2)^T,\quad
 w_r=(0,1,-2r,r^2)^T.
\]
Differentiation of the original polynomial, followed by FRF31, gives
\[
 J_3=a^3v_r+a v_{1,r},\qquad J_4=a^2w_r,\tag{FRF40}
\]
\[
 J_1-\frac{2r}{a^2}J_3=O(1),\qquad
 J_2-\frac2aJ_3-\frac{14ir}{a}J_4=O(1).\tag{FRF41}
\]
These are finite Laurent identities: after the displayed subtractions,
every power of $a$ is nonpositive, with polynomial coefficients in bounded
$A,B,r$. Their elementary column operations and their inverses have bounded
operator norm. Thus all exterior norms have upper bounds of scales
$|a|^3,|a|^5,|a|^5,1$.
The following literal minors give matching lower bounds:
\[
 J_{1,3}=a^3,\quad
 \det J_{\{1,2\},\{3,4\}}=a^5,\quad
 \det J_{\{1,2,3\},(3,4,2)}=a^5(12Ar^2+6r+2B).\tag{FRF42}
\]
The last column list is ordered as written; its sign is retained.
Finally $\det J=-2$. Successive ratios of exterior norms give FRF39.
Fixed positive metrics change their bounding constants, not their exponents.
The executable check expands the full Laurent identities, so no leading
column or concealed minor is omitted.
\end{proof}

For the four colliding-root trajectories in Theorem~\ref{thm:branches},
$r\to1,A,B\to0$ and $F''(r)\to6$. Substitution of FRF37 into FRF39 gives
\[
 \boxed{\sigma(DP_{\mathcal W})\asymp
 (\varepsilon^{-3/4},\varepsilon^{-1/2},1,\varepsilon^{5/4}).}
 \tag{FRF43}
\]
Their squared pullback-metric scales are
$\varepsilon^{-3/2},\varepsilon^{-1},1,\varepsilon^{5/2}$.

For the excluded infinity sign, differentiate FRF38. Since $y_-=O(a)$
along this orbit, the successive leading exterior pole orders of its
transition derivative are $4,5,4,0$. Explicit leading minors have
coefficients $120i,-336i,64i,1$, respectively. These follow from the
entries $2i/a^2,-12i/a^3$, and
\[
 \partial_a w_+=-14iy_-^2/a^2-56y_-/a^3+120i/a^4,\quad
 \partial_{y_-}w_+=28iy_-/a+28/a^2.
\]
The negative-chart derivative has an invertible finite limit.
Using $DP_+=DP_-\,(D C_{-+})^{-1}$ and $a\asymp\varepsilon$ therefore gives
\[
 \boxed{\sigma(DP_{\mathcal W})\asymp
 (\varepsilon^{-4},\varepsilon^{-1},\varepsilon,\varepsilon^4).}
 \tag{FRF44}
\]
The operator condition numbers in FRF43 and FRF44 have scales
$\varepsilon^{-2}$ and $\varepsilon^{-8}$. The corresponding relative
Hermitian Gram condition numbers are their squares,
$\varepsilon^{-4}$ and $\varepsilon^{-16}$. This explicitly distinguishes
the operator and metric conventions. With the same original metric on
both sides of $P_{\mathcal W}$, the product of all four singular values
remains exactly two at every finite point.

There is no unevaluated choice of a four-dimensional metric in these
claims. Let $F=YK^{-1}$ be the matrix of $\Psi$ in target monomials,
$c'=a_{\rm amp}/2-4$, $b_s=s/2$, $M_s=(2\pi)^{s/2}$, where
$s=1$ or $a_{\rm amp}$. The original formula FC56a is
\[
 G=F^*\mathcal H F,\quad
 \mathcal H_{jk}=\sum_{p=0}^j\sum_{q=0}^k
 \binom jp\binom kq(c')^{j+k-p-q}(-i)^p i^q m_{p+q},
 \quad0\le j,k\le3,\tag{FRF45}
\]
\[
 m_0=M_s,\quad m_2=M_sb_s,\quad m_4=M_s(3b_s^2+2b_s),\quad
 m_6=M_s(15b_s^3+30b_s^2+16b_s),\quad m_{2j+1}=0.
\]
These finite formulas evaluate $G_{11},G_{44}$ in FRF35--36 in the original
moments $\mu_v,\ldots,\mu_{v+3}$, quartet, source mass and center.
The full original Gamma norm is used; no Euclidean substitution or
uncomputed affine minimum is hidden in these two constants.

\section{A finite constructive witness for every matrix outside the complete line}
The classification also has a bounded exact witness search; it need not
end with an abstract statement that some trajectory fails.

\begin{proposition}[Finite boundary witness construction, FRF-W]
Let $M$ be any constant matrix outside FRF9.
If the $A$ component of $MY$ has a nonzero $B,C,$ or $D$ coefficient,
one of the $216$ points
\[
 Y_0=(0,B,C,D),\qquad B,C,D\in\{0,1,2,3,4,5\}
\]
satisfies $\Delta(Y_0)(MY_0)_A\ne0$. A small real-time backward segment
of $e^{tM}Y_0$ then supplies the excluded-infinity-sign cubic escape.
If those three coefficients vanish, a transverse simple-double-root
point with $A\ne0$ is obtained from one of the $39^3=59319$ triples
$(A,B,r)\in\{0,\ldots,38\}^3$ using FRF10. It supplies four
quarter-power escapes while four other inverse states stay finite.
\end{proposition}
\begin{proof}
In the first case $\Delta_0(B,C,D)(MY_0)_A$ is a nonzero polynomial of
total degree at most five. A polynomial of degree at most five in each
variable cannot vanish on the whole six-by-six-by-six grid unless it
is zero, by successive univariate interpolation. This proves the first
finite selection. Its leading-coefficient derivative is nonzero and its
quartic remains squarefree near the endpoint, so FRF38 gives the cubic
pole; the other seven signed roots have finite positive-chart limits.

In the second case Theorem~\ref{thm:complete}'s explicit tangency
calculation implies $\mathcal D_M\Delta=(MY)\cdot\nabla\Delta$ does
not vanish identically on the discriminant. Compose it with FRF10 to
obtain a nonzero polynomial $S(A,B,r)$ of total degree at most thirty:
each coordinate has degree at most five, and $\mathcal D_M\Delta$ has
degree at most six. The additional factors
\[
 A,\qquad F''(r)=12Ar^2+6r+2B,\qquad
 D_Q=1-4Ar-8A^2r^2-4AB
\]
have degrees one, three and four. Their product with $S$ is nonzero
and has degree at most thirty-eight. The same interpolation argument
selects a nonzero value on the claimed grid. The resulting quartic is
$(u-r)^2[Au^2+(1+2Ar)u+B+2r+3Ar^2]$; $F''(r)\ne0$ and $D_Q\ne0$
ensure precisely a double root and two other distinct simple roots.
Transversality is $S\ne0$. Taylor expansion as in FRF37 proves the
four quarter-power escapes. The other two roots, with both signs, have
finite limits by FRF7. A sufficiently short real-time backward segment
avoids all other boundaries. Its positive length can be selected by
the same explicit polynomial derivative bounds as FRF29, with the chosen
finite target replacing $(0,0,-3,2)$.
\end{proof}

\section{The unchanged equation (11) and the remaining arithmetic receiver}
At the original finite cutoff $D$, the conductor still has the WCF matrix
\[
 A_{D,s}=\diag(\rho_j^{(s)})^{-1}T_D(g_A)
              \diag(\rho_{j+v}^{(s)}),\quad
 (\rho_n^{(s)})^2=\frac{n!}{(s/2)_n},\tag{FRF46}
\]
\[
 g_A(t)=t^{-v}e^{-4\omega(t)}E_A(\omega(t)),\quad
 \omega(t)=-i\arctan t,\quad
 \det A_{D,s}=\left(\frac{(-i)^v\mu_v}{v!}\right)^{D+1}
              \prod_{j=0}^D\frac{\rho_{j+v}^{(s)}}{\rho_j^{(s)}}\ne0.
\]
All data in FRF46 are fixed during the real-time orbit. The escaping
quantities are inverse branches of $P_{\mathcal W}$, not the entries
of $A_{D,s}$. Its source kernel remains exactly $\mathcal P_{v-1}$.
A change of the conductor's first nonzero moment would be a separate
parameter deformation with the fixed-quotient kernel correction already
specified in the provider's FC25.

For completeness the lifted metric action also remains exact. For a fixed
positive $H$, define $G_P(q)=DP(q)^*HDP(q)$. Differentiating the flow
identity, including both tangent-vector terms, gives
\[
 \mathcal L_{\mathcal X_M}G_P
       =DP^*(M^*H+HM)DP.\tag{FRF47}
\]
Subtracting $2cG_P$ pulls back the actual centered target defect
$M^*H+HM-2cH$; it does not annihilate it. FRF23 and FRF45 return this
calculation to the specified original conductor space. Passing onward
to the programme's observation quotient requires its actual observation,
full Taylor unit, kernel and attained arithmetic Gram. Those are not
replaced by this nonlinear four-dimensional pullback. In particular
neither the remaining signed projected cross product nor the seven
native allocations receives a numerical or asymptotic value from FRF47.

The new proved results are FRF-C (complete-lift classification), FRF-B
(the explicit real-time eight-state census and original-norm constants),
FRF39--44 (the finite-root differential theorem and both mixed-orbit
spectra), and FRF-W (a finite boundary-witness construction). The
source-identity reconciliation is exact at FRF1--5. The original inverse
fibre and monodromy conclusions are inherited from \cite{Signed} and
are not advertised as new again.

\section{Verification and publication scope}
The adjacent \texttt{verify\_real\_flow.py} checks the factor and literal
polynomials against one another, both constraints, the complete Jacobian,
the full sixteen-variable tangency system, the original marked matrix
and its needed inverse columns, the signed transverse coefficients,
the three finite limiting points, the complete finite-root inverse,
its differential columns and exact exterior witnesses, and both charts.
Its optional negative controls reject replacing $KRK^{-1}$ with $R$
and reversing the transverse coefficient. Correctness checks use explicit
exceptions, so optimized Python retains them. The supplied workflow runs
normal and optimized versions and compares their complete receipts.
Actual execution status and run identifiers belong in the separate
verification record; source code is not represented as an executed check.
The arguments about completeness, analytic continuation and limiting
norms are the written proofs above, not consequences of a checker exit code.

\begin{thebibliography}{9}
\bibitem{Signed}
The Clankers. (2026, September 20).
\emph{The complete signed ES--Fable state in the original conductor}.
\texttt{FABLE\_TO\_ORIGINAL\_CONDUCTOR.tex}, complete FC/GF/SE/SM edition.
KokunoYumeto/zeta-function-research-reader, commit
\texttt{d73b545c7e70dd55fb8b58557f704666ba310626}.
Equations FC5--9, FC15a, FC20--35, FC43--44b, FC50--56a and the
complete inverse-fibre sections are the specific inputs used here.
\url{https://github.com/KokunoYumeto/zeta-function-research-reader/tree/d73b545c7e70dd55fb8b58557f704666ba310626/workbenches/splitzero-tandem/continuations/20260920-fable-signed-states}
\bibitem{Tao}
Tao, T. (2026, July 21).
\emph{A digestion of the Jacobian conjecture counterexample}.
What's New. Equations (2), (5)--(10), for the linear--quadratic factor
construction and its source attribution; not for the additional
four-dimensional ES polynomial.
\url{https://terrytao.wordpress.com/2026/07/21/a-digestion-of-the-jacobian-conjecture-counterexample/}
\end{thebibliography}
\end{document}
